% Chapter: The Ladder. Tables under results/elo_duel/ (1v1) and results/elo/
% (free-for-all); case-study figures under results/case_studies/.

\section{From games to ratings}
\label{sec:elo-protocol}

Before any table: what exactly are we measuring, and how do games become
numbers?

\paragraph{The protocols.} All thirteen models of Table~\ref{tab:lineage}
compete. Two tournaments were run with the same harness, seeds, and size
(312 games each; four per pairing, sides alternated, deterministic seeds):

\begin{itemize}
  \item \textbf{Duel} (headline): symmetric 1v1 on a fresh random map;
    winner by elimination, or by node count at the 5000-tick limit. The
    standard basis for a rating.
  \item \textbf{Free-for-all} (secondary): three players --- the rated bot
    in one seat, two \emph{independent, uncoordinated} copies of its
    opponent in the others; the rated bot scores against the
    better-performing copy. A deliberately hostile environment; kept
    because it turns out to discriminate where the duel cannot.
\end{itemize}

The atomic quantity throughout is a pairing's \emph{score rate}: the
average score (win $=1$, tie $=\tfrac12$) one model achieves against
another over their games. A score rate of 0.75 means ``wins three games in
four against this opponent.''

\paragraph{From score rates to a rating.} To compress 78 pairings into one
number per bot we fit a Bradley--Terry model \citep{bradleyterry1952}:
assume each bot $i$ has a strength $p_i$ and that $i$ beats $j$ with
probability $p_i/(p_i+p_j)$, then find the strengths that make the observed
results most likely \citep{hunter2004}, reported on the familiar Elo scale
\citep{elo1978}. This differs from chess-style Elo updates only in being
fit to the whole tournament at once: the result does not depend on game
order and is exactly reproducible from the game log. Two reading rules
follow from the fit. First, with four games per pairing, small rating gaps
are noise --- read the ladder in tiers. Second, when a bot loses nearly
everything, its fitted strength runs off toward $-\infty$; the extreme
negative values at the bottom of the table are scale artifacts, not
meaningful magnitudes.

\section{The ranking}
\label{sec:elo-ladder}

\begin{table}[htbp]
\centering
\resultstable{elo_duel/tables/ladder_combined}
\caption{The ladder under the duel protocol, with the free-for-all ranking
alongside (the two agree: Spearman $\rho = \ffaduelspearman$).}
\label{tab:elo}
\end{table}

Both tournaments reach the same structural conclusions: the same podium
$\{\code{v8}, \code{v9}, \code{v12}\}$, the same \code{v11} collapse at an
identical 0.17 score rate, the same basement. Four headlines:

\begin{enumerate}
  \item \textbf{The duel cannot separate the podium.} \code{v8}, \code{v9},
    and \code{v12} finish in an exact three-way tie (0.90 apiece): each
    crushes everything below and splits the games among themselves. The
    free-for-all \emph{does} separate them (0.90/0.85/0.75) --- being
    outnumbered amplifies small differences. The tie is not just a
    small-sample problem: Chapter~\ref{ch:limits} shows the podium's
    head-to-head structure is genuinely cyclic.
  \item \textbf{\code{v8} shares the top} using the one feature no later
    model kept --- the local softmax (\S\ref{sec:bg-lineage}). The most
    obvious unexplored bot is \code{v8}'s scoring plus \code{v12}'s
    transport.
  \item \textbf{\code{v11} collapsed, in both tournaments.} One feature
    separates it from its parent and one from its child; \S\ref{sec:elo-v11}
    pins down which feature and why.
  \item \textbf{Never attacking is a decent strategy.}
    \code{v0\_expansion}, which has no war agent at all, ranks 4th in duels
    --- above every pre-OT model that fights. The eager attackers
    (\code{v1\_knapsack} and its hybrid) lose both tournaments outright.
\end{enumerate}

\section{The lineage as an ablation study}
\label{sec:elo-ablation}

A flat ranking hides a second, sharper reading of the same games. Because
each lineage step adds (approximately) one feature
(\S\ref{sec:bg-lineage}), the child-vs-parent head-to-head inside the
tournaments is an ablation: the marginal value of one idea with everything
else held fixed. Table~\ref{tab:ablation} reads the existing game logs this
way --- no new games were played.

\begin{table}[htbp]
\centering
\resultstable{elo_duel/tables/ablation}
\caption{Each lineage step as an ablation: the child's head-to-head score
rate against its parent (4 games each per protocol) and the overall rating
gap. $\Delta$Elo magnitudes at the extremes are scale artifacts
(\S\ref{sec:elo-protocol}); the head-to-head columns are the honest
ablation signal.}
\label{tab:ablation}
\end{table}

Four readings:

\begin{enumerate}
  \item \textbf{The biggest single wins are the simplest features.}
    Replacing the literal knapsack with \code{v4}'s priority rule sweeps
    (1.00/1.00), and the border reserve sweeps (1.00/1.00). The most
    sophisticated machinery in the lineage --- the annealed layout ---
    buys \emph{nothing} head-to-head (\code{v4} vs.\ \code{v5\_qubo}:
    0.50/0.50), the gym-vs-game gap of \S\ref{sec:economy} in ablation
    form.
  \item \textbf{Adding war made the bot worse.} The first step,
    \code{v0} $\to$ \code{v1\_knapsack}, loses every game to its pacifist
    parent (0.00/0.00) --- though it is confounded (the step also dropped
    expansion), the aggressive posture never recovers until \code{v4}.
  \item \textbf{The same feature flips sign by chassis.} Construction
    discipline scores 0.75 when added to \code{v6}, and 0/8 when re-added
    to the OT chassis (\code{v9} $\to$ \code{v10}) --- feature value is
    contextual, which is exactly why the ablation column and the global
    rating disagree in places.
  \item \textbf{The \code{v11} step is the largest regression in the
    lineage} (0.00/0.00 against \code{v10}), and the step after it the
    largest recovery. The case study below takes it apart.
\end{enumerate}

\section{Case study: the one-feature collapse of \code{v11}}
\label{sec:elo-v11}

The ladder's strangest datum: \code{v11} sits hundreds of Elo below both its
parent \code{v10} and its child \code{v12}. The lineage
(\S\ref{sec:bg-lineage}) makes this a controlled puzzle, because each step
is a single feature:

\begin{itemize}
  \item \code{v10} $\to$ \code{v11}: nodes stop requesting troops once help
    is on the way (troops en route count as delivered).
  \item \code{v11} $\to$ \code{v12}: border nodes additionally keep an
    always-on request --- the border reserve.
\end{itemize}

\paragraph{The experiment.} Games are deterministic given a seed, so we can
re-run a real tournament game and change nothing but the bot in one seat.
We take FFA game 1301 --- a map where \code{v11}, facing two \code{v10}s,
was destroyed --- and play it three times, once per lineage step, always in
the same seat against the same opponents. Figure~\ref{fig:v11} shows one
game per row; in every row, blue is the bot under test and green the two
\code{v10} opponents.

\paragraph{What happens.}
\begin{itemize}
  \item \textbf{\code{v10} in the seat (top row):} survives to the time
    limit holding about a quarter of the map. The seat is playable.
  \item \textbf{\code{v11} in the seat (middle row):} wiped out by tick
    1000. Its curve is indistinguishable from the \emph{losing} \code{v10}
    opponent's --- it doesn't malfunction visibly, it just loses.
  \item \textbf{\code{v12} in the seat (bottom row):} owns the entire map by
    tick 1000, both opponents dead.
\end{itemize}

Same map, same opponents, one feature per step: switching on
``stop requesting once help is on the way'' turns a playable position into a
loss, and adding the border reserve turns the same position into total
domination. The tournament aggregates match (score rates 0.65 / 0.17 /
0.75).

\paragraph{Why would such an innocuous feature be catastrophic?} Our working
hypothesis: a node under attack needs a \emph{continuous stream} of
reinforcements, but the moment the first ones are dispatched, its request
counts as filled and it goes silent --- so contested nodes stop calling for
help precisely while they are being ground down. The border reserve never
goes silent, which is why \code{v12} recovers. This fits the replay, the
aggregate, and the fact that the collapse is identical under both
tournament protocols --- but it is a hypothesis, not yet an observation.
The direct check (log each node's requests during a replayed loss) is
queued in Chapter~\ref{ch:findings}. Keep the border reserve in mind: this
is the first of three places it proves decisive
(\S\ref{sec:elo-rps}, \S\ref{sec:combat} are the others).

\begin{figure}[p]
\centering
\includegraphics[width=\textwidth]{\resultsfig{case_studies/figures/v10_baseline.pdf}}\\[2pt]
\includegraphics[width=\textwidth]{\resultsfig{case_studies/figures/v11_stall.pdf}}\\[2pt]
\includegraphics[width=\textwidth]{\resultsfig{case_studies/figures/v12_counterfactual.pdf}}
\caption{One tournament game (seed 1301), played three times; only the bot
in the test seat (blue) changes between rows. Top: \code{v10} --- survives
to the time limit. Middle: \code{v11} --- destroyed by tick 1000. Bottom:
\code{v12} --- takes the whole map by tick 1000. Left panels: total troops;
right panels: nodes held.}
\label{fig:v11}
\end{figure}

\bigskip
\noindent
Where this leaves us: a defensible ranking, a feature-by-feature account of
how the bots got here, and one anomaly explained. But the ranking's own
summit is a three-way tie that more games would not break --- \code{v9}
beats \code{v8}, \code{v12} beats \code{v9}, and \code{v12} only splits
with \code{v8}. That is not noise; it is structure a single number cannot
hold. Measuring how much of the game lives beyond any ranking is the next
chapter.
